\documentclass[12pt]{article}
\usepackage{amsmath,amsthm,amssymb,amsfonts,setspace}

\textwidth 7.0 truein
\oddsidemargin -0.25in   %left-hand edge
\evensidemargin -0.5 truein  %right-hand edge
\topmargin -0.85in      %top of paper to top of head, pulls whole unit
\textheight 9.5in

%%%% In most cases you won't need to edit anything above this line %%%%

\begin{document}
\hfill Name 

\hfill Math 2270

\hfill Assignment \#


\begin{itemize}
\item[1)] Question one text

Solution

\item[2)] Question two text

Solution

\item[4)] Define $u_n$ by $u_0 = 0$, $u_1 = 4$, and $u_{n+2} = \frac{6}{5}u_{n+1} - u_n$. Show that $|u_n|\leq 5$ for all $n$.

\[
u_{n+2} - \frac{3+4i}{5}u_{n+1} = \frac{3-4i}{5}(u_{n+1}- \frac{3+4i}{5}u_n)
\]
\[
u_{n+1} - \frac{3+4i}{5}u_{n} = \left(\frac{3-4i}{5}\right)^n(u_1- \frac{3+4i}{5}u_0) = 4\left(\frac{3-4i}{5}\right)^n
\]
\[
u_{n+1} + \frac{5}{2i} \left(\frac{3-4i}{5}\right)^{n+1} = \frac{3+4i}{5}\left(u_{n} +\frac{5}{2i}\left(\frac{3-4i}{5}\right)^n\right)
\]
\[
u_n + \frac{5}{2i} \left(\frac{3-4i}{5}\right)^n=\left(\frac{3+4i}{5}\right)^n\left(u_0+ \frac{5}{2i}\right)
\]
\[
u_n = \frac{5}{2i}\left(\frac{3+4i}{5}\right)^n-\frac{5}{2i} \left(\frac{3-4i}{5}\right)^n
\]

\item[5)] Show that the next integer above $(\sqrt{3}+1)^{2n}$ is divisible by $2^{n+1}$

\[
\begin{aligned}
    (\sqrt{3}+1)^{2n}+(\sqrt{3}-1)^{2n}
    &=\left((\sqrt{3}+1)^{n}+(\sqrt{3}-1)^{n}\right)^2-2(\sqrt{3}+1)^{n}(\sqrt{3}-1)^{n}\\
    &=\left((\sqrt{3}+1)^{n}+(\sqrt{3}-1)^{n}\right)^2-2^{n+1}
\end{aligned}
\]

\item[9)] If $u_0 = 0$, $u_1 = 1$, and $u_{n+2} = 4(u_{n+1} - u_n)$, find $u_{16}$.

\[
u_{n+2} - 2u_{n+1} = 2(u_{n+1} - 2u_n).
\]
\[
u_{n+1} - 2u_n =(u_{1} - 2u_0) \cdot 2^n=2^n
\]
\[
\frac{u_{n+1}}{2^{n+1}} - \frac{u_n}{2^n} = \frac{1}{2}.
\]
\[
\frac{u_n}{2^n} = u_0 + \frac{n}{2}=\frac{n}{2}
\]
\[
u_n = n \cdot 2^{n-1}
\]
So, $u_{16} = 16 \cdot 2^{15} = 2^4 \cdot 2^{15} = 2^{19}$.


\item[10)] Let $a_0 = 1$, $a_1 = 2$, and $a_n = 4a_{n-1} - a_{n-2}$ for $n \ge 2$.
Find an odd prime factor of
$a_{2015}$.

\[
a_{n+1} - (2 + \sqrt{3}) a_n
= (2 - \sqrt{3})(a_n - (2 + \sqrt{3})a_{n-1})
\]
\[
a_{n+1} - (2 + \sqrt{3}) a_n
= (2 - \sqrt{3})^n(a_1 - (2 + \sqrt{3})a_0)=\sqrt{3}(2 - \sqrt{3})^n
\]
\[
a_{n+1}+\frac{(2-\sqrt{3})^{n+1}}{2}=(2+\sqrt{3})\left(a_{n}+\frac{(2-\sqrt{3})^n}{2}\right)
\]
\[
a_{n}+\frac{(2-\sqrt{3})^n}{2}=(2+\sqrt{3})^n\left(a_0+\frac{1}{2}\right)=\frac{3}{2}(2+\sqrt{3})^n
\]
\[
a_{n}=\frac{3}{2}(2+\sqrt{3})^n-\frac{1}{2}(2-\sqrt{3})^n
\]


\item[20 a)] Let $a_0,\dots, a_{k-1}$ be real numbers, and define
\begin{equation}
a_n = \frac{1}{k}(a_{n-1} + a_{n-2} +\dots+ a_{n-k}), \quad n \geq k.
\end{equation}
Find $\lim_{n\rightarrow \infty} a_n$ (in terms of $a_0, a_1,\dots,a_{k-1}$).

\[
L= \frac{1}{k}(a_{0} + a_{1} +\dots+ a_{k})
\]
Define $b_n=a_n-L$, then it has the same property as (1) and also $b_{0} + b_{1} +\dots+ b_{k}=0$. 

Then $\lim_{n\rightarrow \infty} b_n=0$ because every term is 0, then $\lim_{n\rightarrow \infty} a_n=\lim_{n\rightarrow \infty} b_n+L=L$.

\end{itemize}


\end{document}


